% !TEX program = pdflatex
% ============================================================
%  FICHIER UNIQUE POUR OVERLEAF - Si main.tex echoue, compilez CELUI-CI
%  Menu Overleaf : Main document = main-overleaf.tex
% ============================================================
\documentclass[12pt,a4paper,oneside]{report}
\usepackage[T1]{fontenc}
\usepackage[utf8]{inputenc}
\usepackage{lmodern}
\usepackage[french]{babel}
\usepackage[margin=2.5cm]{geometry}
\usepackage{graphicx}
\usepackage{tikz}
\usetikzlibrary{arrows.meta,positioning}
\usepackage{xcolor}
\usepackage{booktabs}
\usepackage{float}
\usepackage{hyperref}
\usepackage{enumitem}
\usepackage{listings}

\definecolor{SultanTeal}{HTML}{0D4F4F}
\definecolor{SultanSand}{HTML}{F5F0E8}

\hypersetup{colorlinks=true,linkcolor=SultanTeal}

\begin{document}

\begin{titlepage}
\centering
\vspace*{2cm}
{\Huge\bfseries\color{SultanTeal} Sultan Morocco\par}
\vspace{1cm}
{\Large Rapport technique -- Application web Laravel 13\par}
\vspace{2cm}
{\large PHP 8.3 / Tailwind v4 / Vite 8\par}
\vfill
{\large Mai 2026\par}
\end{titlepage}

\tableofcontents
\cleardoublepage

\chapter{Introduction}
\textbf{Sultan Morocco} est une application de decouverte du Maroc : hotels et restaurants (donnees statiques PHP), carte Leaflet, chatbot, authentification et messagerie directe.

\chapter{Architecture}
\begin{figure}[H]
\centering
\begin{tikzpicture}[
  box/.style={draw=SultanTeal, rounded corners, fill=SultanSand, minimum width=3cm,
              minimum height=1cm, align=center, font=\small}
]
  \node[box] (b) {Navigateur};
  \node[box, below=1.2cm of b] (l) {Laravel};
  \node[box, below left=1cm and 0.5cm of l] (s) {Listings\\statiques};
  \node[box, below right=1cm and 0.5cm of l] (d) {BDD\\users + msg};
  \draw[-{Stealth}, thick] (b)--(l) (l)--(s) (l)--(d);
\end{tikzpicture}
\caption{Architecture trois couches}
\end{figure}

\chapter{Fonctionnalites}
\section{Hotels (30 entrees)}
\texttt{HotelController::\$HOTELS} sert \texttt{GET /api/hotels} (pagination, tri best\_value ou popularity). Adresses via Nominatim dans le navigateur.

\section{Restaurants (15 entrees)}
\texttt{RestaurantController::\$RESTAURANTS} -- \texttt{GET /api/restaurants}.

\section{Carte}
Leaflet + marqueurs depuis l'API hotels.

\section{Chatbot}
\texttt{POST /api/chat} -- regles par mots-cles, OpenAI optionnel (\texttt{OPENAI\_API\_KEY}).

\section{Messagerie}
Tables \texttt{conversations}, \texttt{conversation\_user}, \texttt{messages}.

\chapter{Routes principales}
\begin{tabular}{@{}lll@{}}
\toprule
Methode & URI & Role \\
\midrule
GET & / & Accueil \\
GET & /hotels & Page hotels \\
GET & /api/hotels & JSON hotels \\
GET & /restaurants & Restaurants \\
GET & /map & Carte \\
POST & /api/chat & Chatbot \\
GET/POST & /login, /signup & Auth \\
GET & /messages & Inbox \\
POST & /messages/\{user\} & Envoi message \\
\bottomrule
\end{tabular}

\chapter{Installation}
\begin{lstlisting}
cd Sultan_Morocco
composer run setup
php artisan serve
\end{lstlisting}
Ouvrir \url{http://127.0.0.1:8000}. Aucune cle API requise pour les listings.

\chapter{Conclusion}
Projet operationnel avec schema DB etendu (places, avis, itineraires) pour evolution future.

\end{document}
